\documentclass[aspectratio=169,10pt]{beamer}
\usepackage[T1]{fontenc}
\usepackage{lmodern}
\usepackage{listings}
\usepackage{xcolor}

\definecolor{Ink}{HTML}{172B3A}
\definecolor{Blue}{HTML}{235789}
\definecolor{Muted}{HTML}{526574}
\definecolor{Paper}{HTML}{F5F7F8}
\definecolor{Code}{HTML}{EAF0F4}
\setbeamercolor{normal text}{fg=Ink,bg=white}
\setbeamercolor{frametitle}{fg=Ink,bg=white}
\setbeamercolor{title}{fg=Ink}
\setbeamercolor{structure}{fg=Blue}
\setbeamercolor{block title}{fg=Blue,bg=white}
\setbeamercolor{block body}{fg=Ink,bg=Paper}
\setbeamertemplate{navigation symbols}{}
\setbeamertemplate{footline}[frame number]
\setbeamersize{text margin left=0.7cm,text margin right=0.7cm}
\setlength{\parskip}{0.35em}
\lstset{basicstyle=\ttfamily\small,columns=fullflexible,keepspaces=true,
  showstringspaces=false,breaklines=true,backgroundcolor=\color{Code},
  frame=single,rulecolor=\color{Code},xleftmargin=0.4em,xrightmargin=0.4em,
  aboveskip=0.6em,belowskip=0.6em}

\title{Lazy Evaluation with Streams}
\subtitle{Delaying work until a value is needed}
\author{CS413}
\date{Fall 2026}

\begin{document}
\begin{frame}[plain]\titlepage\end{frame}

% 2 minutes
\begin{frame}[fragile]{Evaluation can happen now or later}
  In a strict language, a function's arguments are evaluated before the call.
  Lazy evaluation delays a computation until its result is requested.
  \begin{columns}[T]
    \column{.48\textwidth}
      \textbf{Strict argument}
      \begin{lstlisting}
choose(42, expensive())
      \end{lstlisting}
      Both arguments are evaluated before \texttt{choose} runs.
    \column{.48\textwidth}
      \textbf{Delayed argument}
      \begin{lstlisting}
choose(42, thunk(expensive))
      \end{lstlisting}
      The second computation can wait until it is needed.
  \end{columns}
  \vspace{0.5em}
  This lecture uses explicit thunks and lazy streams, so the delay is visible
  in the program rather than built into every expression.
\end{frame}

% 2 minutes
\begin{frame}[fragile]{A thunk packages a deferred computation}
  A thunk is a function with no meaningful input that computes a value later.
  \begin{lstlisting}
Lazy[T] = () -> T
  \end{lstlisting}
  In LAMBDA, the interface is written using \texttt{lazy\_make} and
  \texttt{lazy\_eval}:
  \begin{lstlisting}
deferred : Lazy[sint] =
  lazy_make(lam (u: unit) => 6 * 7)
answer = lazy_eval(deferred)
  \end{lstlisting}
  Creating \texttt{deferred} stores the computation; \texttt{lazy\_eval}
  requests its value. In this runtime, a lazy value caches its first result.
\end{frame}

% 2 minutes
\begin{frame}[fragile]{A lazy stream delays the next cell}
  A finite list stores its whole spine. A stream can store one cell now and
  delay the rest.
  \begin{lstlisting}
Stcn[T] = stcn_nil()
       | stcn_cons(head: T, tail: Strm[T])
Strm[T] = Lazy[Stcn[T]]
  \end{lstlisting}
  \begin{itemize}
    \item \texttt{stcn\_nil} marks the end.
    \item \texttt{stcn\_cons(x, xs)} stores a head and a delayed stream tail.
    \item The outer \texttt{Lazy} also delays inspection of the current cell.
  \end{itemize}
  This is the course library's explicit representation of a lazy list.
\end{frame}

% 3 minutes
\begin{frame}[fragile]{Constructing a stream delays its constructor}
  The library wraps a stream cell in a thunk:
  \begin{lstlisting}
strm_nil[T] : unit -> Strm[T] = lam (u) =>
  lazy_make(lam (v) => stcn_nil())

strm_cons[T] : (T, Strm[T]) -> Strm[T] = lam (p) =>
  lazy_make(lam (u) => stcn_cons(p.0, p.1))
  \end{lstlisting}
  \begin{itemize}
    \item Calling \texttt{strm\_cons} creates a stream value without inspecting its cell.
    \item \texttt{strm\_head(xs)} forces the outer thunk, then reads the head.
    \item \texttt{strm\_tail(xs)} gives the tail stream; its cell stays delayed.
  \end{itemize}
  Operations such as \texttt{strm\_head} need a nonempty stream.
\end{frame}

% 3 minutes
\begin{frame}[fragile]{Demand propagates only as far as the consumer asks}
  Suppose the stream has cells \texttt{1, 2, 3, ...}. To get the first head:
  \begin{enumerate}
    \item Force the outer lazy value to reveal a \texttt{stcn\_cons}.
    \item Read its head, \texttt{1}.
    \item Leave its tail untouched.
  \end{enumerate}
  \begin{lstlisting}
strm_head(xs)               # forces one cell
strm_head(strm_tail(xs))    # forces two cells
  \end{lstlisting}
  A consumer controls how much of a stream gets evaluated. Reusing a forced
  lazy cell reuses its cached result in this implementation.
\end{frame}

% 3 minutes
\begin{frame}[fragile]{A stream can describe an unbounded sequence}
  Define a stream of integers starting at \texttt{n}; the recursive call sits
  inside a thunk, so it does not run while the current cell is built:
  \begin{lstlisting}
naturals_from : sint -> Strm[sint] = fix from(n) =>
  strm_cons(n,
    lazy_make(lam (u: unit) => lazy_eval(from(n + 1))))
  \end{lstlisting}
  \begin{itemize}
    \item Constructing \texttt{naturals\_from(0)} returns after building its first cell.
    \item The tail computation builds the next cell only when demanded.
    \item The stream has no final cell, but any finite prefix can be observed.
  \end{itemize}
  \texttt{strm\_range(lo, hi)} in the library is the corresponding finite example.
\end{frame}

% 3 minutes
\begin{frame}[fragile]{Mapping transforms cells as they are requested}
  \texttt{strm\_map} delays each output cell and maps the head when that cell
  is forced:
  \begin{lstlisting}
strm_map(xs, lam (x) => x * 2)
  \end{lstlisting}
  For a source stream \texttt{1, 2, 3, ...}, its result describes
  \texttt{2, 4, 6, ...}. Creating the mapped stream does not traverse the source.
  \begin{lstlisting}
strm_to_list(
  strm_take(strm_map(naturals_from(1), lam (x) => x * 2), 4))
# [2, 4, 6, 8]
  \end{lstlisting}
  \texttt{take} makes the finite demand explicit before conversion to a list.
\end{frame}

% 3 minutes
\begin{frame}[fragile]{Filtering may search ahead for the next output}
  \texttt{strm\_filter(xs, p)} skips input cells until it finds one satisfying \texttt{p}.
  \begin{lstlisting}
strm_filter(naturals_from(0), lam (n) => n > 5)
  \end{lstlisting}
  Asking for the first output forces inputs \texttt{0} through \texttt{6}.
  Asking only whether an earlier cell matches can stop sooner:
  \begin{lstlisting}
strm_exists(xs, lam (n) => n == 3)
  \end{lstlisting}
  \texttt{strm\_exists} returns as soon as it finds a match. If the predicate
  never matches an infinite stream, finding the first filtered cell never finishes.
\end{frame}

% 3 minutes
\begin{frame}[fragile]{The sieve generates primes on demand}
  Start with the stream \texttt{2, 3, 4, 5, ...}. Keep its head as the next
  prime, filter multiples of that prime from the tail, then sieve the result:
  \begin{lstlisting}
sieve(xs) =
  let p = strm_head(xs)
  in strm_cons(p,
       lazy_make(lam (u: unit) =>
         lazy_eval(sieve(strm_filter(
           strm_tail(xs), lam (n) => n mod p != 0)))))

primes = sieve(naturals_from(2))
  \end{lstlisting}
  \begin{itemize}
    \item The first output is \texttt{2}; its delayed tail removes multiples of 2.
    \item The next head is \texttt{3}; the next delayed tail removes multiples of 3.
    \item Taking a finite prefix forces only enough filtering and sieving to produce it.
  \end{itemize}
  \texttt{strm\_filter} here denotes the library filter; the recursive call is
  delayed inside the tail thunk.
\end{frame}

% 3 minutes
\begin{frame}{Lazy composition avoids intermediate lists}
  A pipeline can describe a sequence of transformations:
  \begin{center}
    source stream \quad $\longrightarrow$ \quad map \quad $\longrightarrow$
    \quad filter \quad $\longrightarrow$ \quad take
  \end{center}
  \begin{itemize}
    \item Each stage produces a stream whose later cells remain delayed.
    \item The consumer's request pulls values through the pipeline as needed.
    \item No full intermediate list is required.
  \end{itemize}
  Laziness can save work and support infinite sources. It does not promise that
  every pipeline uses less memory or runs faster; retained thunks can keep
  earlier data reachable.
\end{frame}

% 2 minutes
\begin{frame}[fragile]{Short-circuit consumers need only a prefix}
  \texttt{strm\_exists} stops at the first matching item; \texttt{strm\_forall}
  stops at the first item that fails the predicate.
  \begin{lstlisting}
strm_exists(xs, lam (x) => x == 1)
  \end{lstlisting}
  If the first cell of \texttt{xs} is 1, the result is \texttt{true}; the
  consumer need not inspect any later cells. A delayed tail could diverge or
  raise an error, and this result would still be available.
  \begin{block}{Key idea}
    A lazy program can avoid computations that its result does not depend on.
  \end{block}
\end{frame}

% 2 minutes
\begin{frame}{Laziness has costs and limits}
  \begin{itemize}
    \item A thunk needs a closure and may retain the values it captures.
    \item Memoization avoids recomputing a forced cell, but cached values occupy memory.
    \item A consumer that asks for the whole stream must still do all that work.
    \item A nonproductive definition may loop before producing its first cell.
  \end{itemize}
  Ask two questions when evaluating a lazy pipeline: which cells are forced,
  and which values remain reachable afterward?
\end{frame}

% 2 minutes
\begin{frame}{Summary}
  \begin{itemize}
    \item Lazy evaluation delays work until a consumer demands a value.
    \item A thunk represents a deferred computation; this LAMBDA runtime memoizes its result.
    \item A stream is a lazy cell whose tail is another stream.
    \item Mapping and filtering can operate on unbounded inputs, while consumers such as \texttt{take} and \texttt{exists} bound demand.
  \end{itemize}
  \vspace{0.7em}
  \begin{block}{Quick check}
    In \texttt{strm\_exists(xs, p)}, why might it be incorrect to force the
    tail before testing the head?
  \end{block}
\end{frame}
\end{document}
